% page 377 du fac-simile (image p-376 du scan)
% bloc 162.6 x 242.0 mm ; marge gauche 39.4 mm
\pgc{17.780mm}{0.000mm}{
\l{\cel{54}369}
\l{}
\l{metamorfosar. (\sou{trans.}) Igar (ento) pasar de lua formo naturala}
\l{\cel{3}a formo altra. - DEFIRS.}
\l{}
\l{metano. (\sou{kemio)}. Termino unesma CH4 ek la hidokarbidi saturita.}
\l{\cel{3}- DEFIS.}
\l{}
\l{metastabila. (\sou{fiziko)}. Dicesas pri stando fizikala partikulara -}
\l{\cel{2}o di kristalo, o di mixuro - obtenita en kondicioni specala}
\l{\cel{3}di temperaturo : stando nestabila teoriale, ma qua povas}
\l{\cel{3}tamen persistar dum periodo tre longa, o mem senfine.- DEF.}
\l{}
\l{metatarso. (\sou{anat.)}\cel{2}Unionuro de la kin osti paralela ye qua}
\l{\cel{3}konsistas la parto di la pedo qua esas inter tarso e fingri.}
\l{\cel{3}- EFIS.}
\l{}
\l{\sou{metatezo}. (\sou{gram.)}\cel{2}Permuto di literi en vorto. (ex.:\cel{1}\sou{prato}}
\l{\cel{3}vice\cel{1}\sou{parto}). - DEFIRS.}
\l{}
\l{metempsikoso.\cel{2}Doktrino qua agnoskas existi intersequanta,}
\l{\cel{3}ye qui la anmo pasas, de la korpo quan lu vivigis, aden}
\l{\cel{3}olta di altra ento. - DEFIRS.}
\l{}
\l{\sou{meteoro}. (\sou{fiziko}). Fenomeno qua eventas en atmosfero : ciel-}
\l{\cel{3}arko,pol-auroro, pluvo, roso, nivo, vento, trombo, e c.-}
\l{\cel{3}DEFIRS.}
\l{}
\l{\sou{meteorologio.}\cel{1}Ta parto di fiziko qua traktas la fenomeni qui}
\l{\cel{3}eventas en atmosfero. - DEFIRS.}
\l{}
\l{\sou{metilo}. (\sou{kemio}). Radiko monovalenta CH3, termino unesma di}
\l{\cel{3}la serio de la radiki di la karbidi grasoza saturita, e qua}
\l{\cel{3}derivesas de metil-alkoholo per supreso di oxihidrido. -}
\l{\cel{3}DEFIS.}
\l{}
\l{\sou{metileno}. (\sou{kemio)}. Radiko divalenta CH2, derivita de metano}
\l{\cel{3}per la supreso di du atomi di hido. - DEFIS.}
\l{}
\l{\sou{metodo}. I. Iro-modo segun-raciona quan onu praktikas por atin-}
\l{\cel{3}gar sua skopo. - II. Iro-modo segun-raciona quan praktikas}
\l{\cel{3}la spirito por atingar lo exakta, o por demonstrar ta exak-}
\l{\cel{3}tajo pos trovir lu. - DEFIRS.}
\l{}
\l{\sou{metodologio}. Ta parto di logiko qua studias \textquotedbl{}a posteriori\textquotedbl{} la}
\l{\cel{3}metodi di la sorti diversa di konoci, di savi, e plu parti-}
\l{\cel{3}kulare la metodi di la cienci diversa. - DEFIS.}
\l{}
\l{\sou{metonimio.}\cel{1}(\sou{gram.)}\cel{2}Figuro di retoriko, uzata por expresar}
\l{\cel{3}kozo per termino qua indikas altra kozo qua esas unionita}
\l{\cel{3}a lu per relati necesa (kauzo po efekto, kontenanto po}
\l{\cel{3}kontenato, e c.) - DEFIRS.}
\l{}
\l{metonomazio. Transformeso subisata da persono-nomo quan onu}
\l{\cel{3}transportas aden linguo stranjera per tradukar lu (Dubois =}
\l{\cel{3}Sylvius ; Schwarzerd =\cel{1}\sou{Melanchton}\cel{1}; Pesch =\cel{1}\sou{Persiko)}. DEFIRS.}
}
